\documentclass[11pt]{article}
\usepackage[T1]{fontenc}
\usepackage{lmodern,amsmath,amssymb,amsthm,mathtools}
\usepackage[margin=1in]{geometry}
\usepackage[hidelinks]{hyperref}
\emergencystretch=1em
\newtheorem{theorem}{Theorem}\newtheorem{lemma}[theorem]{Lemma}
\newcommand{\R}{\mathbb R}\newcommand{\Z}{\mathbb Z}
\newcommand{\Q}{\mathbb Q}\newcommand{\C}{\mathbb C}
\newcommand{\T}{\mathbb T}\newcommand{\N}{\mathbb N}
\newcommand{\ee}{\mathrm e}\newcommand{\DR}{\mathrm{DR}}
\newcommand{\Hor}{\operatorname{Horiz}}\newcommand{\Tay}{\operatorname{Taylor}}
\title{Sections 9--11 at every rank step}
\author{Author proposal: /root/astra\_v14\_baseline\_gaussian}
\date{7 November 2026}
\begin{document}\maketitle

\noindent\textbf{Status.} This is a source-to-consumer project reconstruction,
not independent acceptance. It separates two analytic supplier contracts,
\textup{F} and \textup{D}, from the algebraic, selection and representation
joins proved here. Their separate reconstruction and review remain required.
The universal group uses the explicit fixed-rank weighted presentation below.
No final inverse theorem is used as a premise.

\section{The family at an arbitrary rank}
Fix a standard $s\geq3$, $m=s-1$ and $1\leq R\leq m$, and an unbounded
limit integer $N$. All groups, lattices, metrics and finite dimensions below
are standard unless a natural ultrapower is displayed or understood.
Braces use $(-1/2,1/2]$ throughout this packet. The Section 12 constructor
allows either fixed half-open fundamental cube; use this same choice there.
Let $H\subseteq[-N,N]\cap{}^*\Z$ be dense and internal. The input consists of
a bounded function $f$ on $[N]$, a global unit-vector multidegree $(1,m)$
nilcharacter $A(h,n)$, and an internal family of unit-vector degree--rank
$(m,R)$ nilcharacters $\chi_h$. The latter have a fixed finite list of rational
models and standard uniform Lipschitz bounds. For every $h\in H$,
$\Delta_h f$ has uniformly positive degree-at-most-$m-1$ correlation with
$A(h,\cdot)\otimes\chi_h$, with bounded internal witnessing tests.
Every vertical identity is on the full natural extension of its group and
compact quotient, not merely their standard points.

This is exactly the new input in rank-iteration001, Section 3. In particular
no core/petal representation or independence statement is inherited from
the preceding rank. They will be reconstructed. The corrected LP03 gives,
after any dense internal restriction of $H$, many additive quadruples
$h_1+h_2=h_3+h_4$ for which
\begin{equation}\label{eq:quad}
 \chi_{h_1}(n)\otimes\chi_{h_2}(n+h_1-h_4)
 \otimes\overline{\chi_{h_3}(n)}
 \otimes\overline{\chi_{h_4}(n+h_1-h_4)}
\end{equation}
is $(m-1)$-biased, with standard uniform witnesses. This hereditary statement
means reapplying LP03 to the restricted correlation hypothesis. It does not
assert that a previously chosen dense set of quadruples survives restriction.

\paragraph{Analytic contract F.}
For internal functions $X_1,\ldots,X_4:H\to{}^*\T$ and $\epsilon=N^{-j}$,
if $X_1(h_1)+\cdots+X_4(h_4)=O(\epsilon)$ on many additive quadruples,
there are a dense internal $H'$, a standard finite $K$, and internal
$\alpha_k\in{}^*\T$, $\beta_k\in{}^*\R$, $\delta\in{}^*\T$ such that
\[
 X_4(h)=\sum_{k=1}^K\{\alpha_k h\}\beta_k+\delta+O(\epsilon)
 \quad(h\in H').
\]
The errors have a common standard bound. A permutation of the four positions
changes which function is recovered. This is the precise Appendix F.1 input;
the signed-graph and no-wrap argument is a separate analytic supplier.

\paragraph{Analytic contract D.}
For a polynomial sequence in one variable in a fixed rational filtered
nilmanifold there is a factorization $g=\varepsilon g'\gamma$ with
$\varepsilon$ bounded, Lipschitz $O(1/N)$ and positive-index Taylor
coefficients $O(N^{-i})$; $g'$ takes values in a connected rational subgroup
$Q$ and is totally equidistributed on every standard finite cover of
$Q/(Q\cap\Gamma)$; and $\gamma$ takes values in a fixed discrete overgroup
$\widehat\Gamma\supseteq\Gamma$ and is periodic modulo $\Gamma$ with standard
period. The subgroup $Q$ is allowed to equal $G$. All factors remain adapted
to the ordinary degree filtration. Only this one-variable contract is needed.
Its quantitative published supplier and corrected proof are separate.

\begin{theorem}\label{thm:output}
Subject to these two analytic contracts and the stated LP03 input, a dense
internal subfamily has the following representation. There are fixed standard
dimension blocks $C,L,B$, fixed filtered lattice-compatible maps from a
weighted universal group $U$, and one integral top character $\eta$ on
$U_{(m,R)}$. Up to a fixed coordinate isometry $\iota$,
\[
 \iota\chi_h(n)=F_h(g_h(n)\Gamma,g_{0,h}(n)\Gamma_0)
 \qquad(n\in[N]).
\]
The lower model has degree--rank at most $(m,R-1)$. The functions $F_h$ are
unit vectors with uniform standard bounds and exact top-character law;
they need not be equal as functions. The horizontal coefficients are fixed
core frequencies, linear petals $\{\alpha_a h\}\beta_a$ for $a\in L$, and
independent petals for $a\in B$. The linear petals are Freiman homomorphisms
on the retained set. For almost every additive quadruple, the core, all four
independent-petal blocks and any three linear-petal blocks are independent
modulo $O(N^{-i})$ separately at each weight $i$.

For every weight-$m$, length-$R$ iterated commutator $w$, the real number
$\eta(w)$ is zero if $w$ contains two petals, or contains one letter from
$B$. These are precisely the annihilations required by LP04. They hold for
every new rank family satisfying the initial hypotheses, without inherited
frequency information. The construction of the final two-variable test is
the separate Section 12 join.
\end{theorem}

\section{Selection and the corrected universal representation}
We use the following elementary internal-choice rule, which also makes the
countability uses precise. For an increasing internal sequence of witness
sets $S_k$, suppose each member of an internal parameter set has a witness
in some standard $S_k$. The least successful index is an internal integer.
It is uniformly bounded: otherwise the internal nonempty sets on which it
exceeds $1,2,\ldots$ have a common element by countable saturation, contrary
to pointwise standard boundedness. At representative indices choose the
first witness in a fixed choice order. This gives internal choices with one
standard bound. Countable saturation here follows by diagonalizing nested
finite intersections: at representative index $v$, choose a witness for the
largest $k\leq v$ whose first $k$ conditions hold there. This $k$ tends to
infinity along the ultrafilter. The unsupported diagonal $x_{v,v}$ alone
would not establish this conclusion.

We also need the ``almost all'' form. For nested internal bad sets $B_k$
in a hyperfinite set $X$, if every $B_k$ has infinitesimal relative size,
the same diagonal construction chooses an unbounded $K$ for which $B_K$
has infinitesimal relative size and contains every standard $B_k$.
Consequently, if their countable union is not contained in a sparse internal
set, some $B_k$ is dense. Applying finite pigeonhole inside that $B_k$ fixes
an integer coefficient vector and a standard error bound. This is the exact
quantifier step used for failed frequency independence below.

First restrict to one of the fixed models for $\chi_h$. Bounded
tautification and normalization use their separately proved project
interfaces. In a fixed model, only finitely many integral vertical characters
can occur under a common Lipschitz bound. Indeed, along a fixed unit-speed
top-torus direction $X$, the identity
$F(\exp(tX)x)=\ee(t\eta(X))F(x)$ and $\|F(x)\|=1$ give
$2\pi|\eta(X)|\leq C_X\operatorname{Lip}(F)$ by taking small $t$.
On an integral basis this bounds all integer character coordinates. Finite
pigeonhole therefore fixes the character.

For each weight $i$, choose finitely many elements of $\Gamma\cap G_{(i,0)}$
whose logarithms generate its rational Lie algebra. Such elements are obtained
by clearing denominators in an adapted rational Mal'cev basis; their real
generated subgroup is $G_{(i,0)}$. Use one formal letter of weight $i$ for
each. Define $U$ by killing all Lie words of weight greater than $m$ and all
weight-$m$ words of length greater than $R$. Set
\[
 U_{(d,r)}=\exp\operatorname{span}\{\text{surviving words of weight}>d,
 \text{ or weight}=d\text{ and length}\geq\max(r,1)\}
\]
for positive $d$, with $U_{(0,0)}=U$. This is the exact LP04 weighted
presentation before the additional petal relations. Its Hall-coordinate
lemma supplies the lattice generated by the letters and its filtered
intersections. The upward-closed word conditions prove every bracket
inclusion. In particular $U_{(d,0)}=U_{(d,1)}$.

Mapping each letter to its chosen lattice element gives a filtered
homomorphism $\phi:U\to G$. The killed words vanish by the degree--rank
bracket law in $G$; the map sends the generated lattice into $\Gamma$.
It is onto each $G_{(i,0)}$. In
$\Hor_i(G)=G_{(i,1)}/G_{(i,2)}$ the images of the weight-$i$ letters span
the real vector space and generate a full sublattice. Thus the induced
map of compact horizontal tori is onto. Select internal preimages of
$\Tay_i(g_h\Gamma)$ to get frequencies $\xi_{h,i,a}\in{}^*\T$.
Pull the fixed top character back by $\phi$ and initially take $G_0$ trivial.
This proves the exact Definition 9.11 representation, including lattice
compatibility and all-limit equivariance, for every new family.

Here and below horizontal Taylor coefficients are homomorphisms on the
polynomial-sequence group, are unchanged by translating $n$, and for
$c\in G_{(i,1)}$ the sequence $c^{\binom ni}$ has only the $i$th horizontal
coefficient $c$ and zero others. These facts use the iterated-difference
commutator rule B.7: all noncommutative correction terms have rank at least
two and disappear modulo $G_{(i,2)}$; an extra difference has weight above
$i$. Binomial differences give the last assertion. Polynomial products use
B.4. No claim that arbitrary Taylor factors commute is used.

If $\xi_{i,a}=\sum_b c_{i,a,b}\xi'_{i,b}$ with standard integers $c_{i,a,b}$,
map the new weight-$i$ letter $e'_{i,b}$ to
$\prod_a e_{i,a}^{c_{i,a,b}}$ in a fixed order. The universal relations and
lattice inclusion hold, giving $\psi:U'\to U$. In the abelian horizontal
quotient this is exactly the indicated integer matrix. Replace $\phi$ by
$\phi\psi$ and $\eta$ by $\eta\psi$; the original orbit and test are unchanged.
When these matrices depend on $h$, the preceding bounded-choice rule and
finite pigeonhole fix them on a dense internal subfamily before this step.

\section{Frequency regularization with the actual quantifiers}
Fix one weight and put $\epsilon=N^{-i}$. A partial description consists of
fixed core frequencies, independent-candidate petals $b_h$, bracket-linear
petals $l_h$, fixed rational frequencies, and $O(\epsilon)$ errors, with all
original frequencies standard integer combinations of them. Every list is
standard finite, and every family is internal. Dense restriction fixes the
integer matrices and error bounds. Put all original frequencies in $b_h$
to obtain an initial partial description.

Minimize lexicographically the numbers of $b$ petals, then $l$ petals,
then core frequencies, among all such partial descriptions on dense
subsets. Bracket-linear petals are required to be individually Freiman
homomorphisms. This costs only a finite dense restriction: place each
$\{\alpha h\}$ in one of a fixed finite collection of intervals of width
less than $1/4$. For an additive quadruple the difference of the four
fractional parts is an integer of absolute value less than one, hence zero.
Multiplication by any real $\beta$ preserves that identity. The intervals
may be anywhere in the fundamental cube; none is required to surround zero.

If independence of the core, four $b$ blocks and the first three $l$ blocks
fails on more than a sparse set of quadruples, the bad-set rule above fixes
one nonzero standard integer relation with error $O(\epsilon)$ on many
quadruples. There are three exhaustive cases.

If it contains only core coefficients, solve for a core with nonzero
coefficient $a$. Choose internal $a$th roots of the other core frequencies.
Division introduces a residue in the finite group $a^{-1}\Z/\Z$ and an
$O(\epsilon)$ error. Move them to the rational and error lists. This reduces
the core count without changing either petal count.

If the relation involves $l$ but no $b$, it is a relation on three free
parameters. Fix two by averaging over the many triples. On a dense set of
the remaining parameter, solve for one $l$ coefficient. Divide the remaining
bracket-linear coefficients as real numbers $\beta/a$, and root the core
coefficients. Freeze the finite rational carry and add the values at the
two fixed parameters to the core. A further fixed-box restriction restores
the Freiman laws for the divided bracket-linear terms. This reduces the
$l$ count without changing the $b$ count.

If a $b$ coefficient occurs, use a symmetry of $h_1+h_2=h_3+h_4$ to place
a nonzero one in position four. The Freiman laws allow elimination of the
fourth $l$ block after this permutation. Group the relation as
$X_1(h_1)+\cdots+X_4(h_4)=O(\epsilon)$, with $X_4$ containing that $b$
coefficient and fixed core terms. Contract F expresses $X_4$ as finitely many
bracket-linear terms plus a constant and error. Solve for the chosen $b$
coefficient, using roots of the other $b$ and core frequencies and dividing
the real bracket-linear coefficients. Freeze the finite rational carry,
enlarge the other lists, and refine into fixed boxes as above. The original
frequencies are still integer combinations of the new lists, because the
old frequencies are integer multiples of their chosen roots. The $b$ count
has decreased. Each case contradicts the lexicographic minimum.

Thus the required independence holds for almost all quadruples. Apply this
argument successively for the finitely many weights $1,\ldots,m$. Dense
restrictions preserve previously obtained independence: their bad sets
remain $o(N^3)$, while the retained parameter set has standard positive
density. Standard integer relations, rational carries, error constants,
finite dimensions and fixed model choices are selected internally with
the explicit bounded-choice rule. This proves the regularization step of
10.5--10.7 conditional only on contract F, with its needed family uniformity.

\section{Removing rational and small coordinates without losing unit norm}
Apply the integer change-of-basis construction to the regularized frequency
lists. We now remove each rational or small coordinate in turn. The
remaining top character and lattice map must be updated; forgetting them
would not give a new representation.

For a small coordinate $\xi=O(N^{-i})$, choose its small real lift and put
$c=\phi(e_{i,a})^{\xi}\in G_{(i,1)}$. Set
$g'_h(n)=c^{-\binom ni}g_h(n)$. B.4 proves polynomiality and horizontal
homomorphism shows that exactly the selected coordinate has been removed.
On $[N]$, $c^{\binom ni}$ is a bounded smooth function of $u=n/(2N)$.
Indeed its logarithm is $\binom{2Nu}{i}\xi\log\phi(e_{i,a})$; its coefficients
and all derivatives on a fixed compact interval are bounded because
$N^i\xi=O(1)$. Extend this logarithm to the unit circle by using its real
formula on the lift $(-1/4,3/4)$ and multiplying by a fixed smooth cutoff
equal to one on $[0,1/2]$ and zero near the lift's endpoints. Let $c_h(u)$
be its exponential. Replace the test by
\[
 F'_h(x,x_0,u)=F_h(c_h(u)x,x_0).
\]
It is globally a unit vector and uniformly Lipschitz. The top image
$\phi(U_{(m,R)})$ is central, so the same exact vertical law holds everywhere.
Add the circle orbit $n/(2N)$ to the lower model. It has degree one, which
is at most $(m,R-1)$ because $m\geq2$. Thus this is an exact representation,
not an approximation to the original sequence.

For a rational coordinate $q\in\Q/\Z$, put $\gamma=\phi(e_{i,a})^q$ and
$g'_h(n)=g_h(n)\gamma^{-\binom ni}$. Again only that horizontal coordinate
is removed. The element $\gamma$ is rational relative to $\Gamma$.
The subgroup $\widehat\Gamma=\langle\Gamma,\gamma\rangle$ is discrete:
in rational Mal'cev coordinates the finite nilpotent collection formulas
express every word as a polynomial in integer exponents of a finite set
of generators and their finitely many commutators; clearing their fixed
rational denominators gives a discrete coordinate grid. Since $\Gamma$
is cocompact, its index in $\widehat\Gamma$ is finite. Its normal core
$\Gamma'=\bigcap_{\delta\in\widehat\Gamma}\delta\Gamma\delta^{-1}$ is
therefore a finite-index subgroup, normalized by $\gamma$.

Choose a standard $b$ with $\gamma^b\in\Gamma$ and a standard period $q_0$
such that $\binom{n+q_0}{i}\equiv\binom ni\pmod b$. For example a sufficiently
divisible multiple of $b\,i!$ works, by the integer binomial difference
formula. Put $\gamma_r=\gamma^{\binom ri}$ for $r\pmod{q_0}$. Then
\[
 F_h(g_h(n)\Gamma,x_0)=F_h(g'_h(n)\gamma_r\Gamma,x_0)
 \quad(n\equiv r\pmod{q_0}).
\]
The test $W_{r,h}(x\Gamma',x_0)=F_h(x\gamma_r\Gamma,x_0)$ is well-defined,
unit, and has the same top-frequency law. Also
$W_{0,h}(x\Gamma',x_0)=F_h(x\Gamma,x_0)$ is such a global unit test.

To interpolate the finitely many residue tests, choose disjoint small arcs
around $r/q_0$ in $\T$. In the $r$th arc choose a smooth angle $\theta_r$
which is zero at $r/q_0$ and $\pi/2$ near its boundary. Set $a_r=\cos\theta_r$
there and zero outside; set $b=\sin\theta_r$ there and $b=1$ outside all
arcs. The supports of the $a_r$ are disjoint and $\sum_r a_r^2+b^2=1$.
The test
\[
 \left(\sum_r a_r(u)W_{r,h}(x,x_0),\ b(u)W_{0,h}(x,x_0)\right)
\]
is globally Lipschitz, unit and has the same vertical character. At the orbit
$u=n/q_0$ it is exactly $(\chi_h(n),0)$. This fixed zero padding is explicit;
we do not assert that an arbitrary scalar interpolation stays unit or that
the old dimension has been retained. Correlation and subsequent matrix
comparisons recover the original family by the fixed coordinate projection.

The map from the remaining universal group must now land in the covering
lattice $\Gamma'$. Choose a standard $M$ for which every remaining lattice
generator's $M$th power lies in $\Gamma'$. Send each new letter to that
$M$th power, and divide its horizontal frequency by $M$. Pull back the
character by the same map. For linear petals choose the division as
$\{\alpha h\}(\beta/M)$; for core coordinates choose fixed roots and for
independent coordinates choose internal roots. Repeating this finite process
produces a common standard denominator. It is not silently dropped, as in
the displayed conclusion of the printed proof of 10.10.

Independence is preserved: multiply an alleged standard integer relation
between divided frequencies by the common denominator to contradict the
old independence. After dividing linear petals, refine again into the fixed
short boxes to make their real fractional-part sums exactly additive, hence
their divided values Freiman. This is valid at arbitrary box locations.
All rational covers, maps, dimensions, and periods are fixed on a dense
internal subfamily before use. Small-coordinate curve bounds are uniform
by the internal bounded-choice rule. A finite composition of paddings is
one fixed coordinate isometry. We have proved the precise representation
part of Theorem~\ref{thm:output}, without replacing the varying $F_h$ by a
single equal function through an approximation argument.

\section{The horizontal consequence of factorization, with denominators}
Here we derive the needed D.6 conclusion from contract D. View a degree--rank
group as ordinarily filtered by $G_i=G_{(i,0)}$. We retain the finer horizontal
spaces $\Hor_i(G)=G_{(i,1)}/G_{(i,2)}$. For a polynomial orbit $g\Gamma$ put
\[
 \Xi_i=\{\lambda\in\widehat{\Hor_i(G/\Gamma)}:
                 \lambda(\Tay_i(g\Gamma))=O(N^{-i})\pmod1\}.
\]
Characters here have standard integer coordinates. Their real linear lifts
define $\Xi_i^\perp$ inside the real horizontal vector space.

Suppose $F$ is a bounded Lipschitz test with nonzero standard average.
Take one scalar component if needed, and factor $g=\varepsilon g'\gamma$.
On a progression of the standard periodic modulus, choose a bounded standard
representative $\gamma_P\in\widehat\Gamma$ of the constant coset of $\gamma$.
Set $Q_P=\gamma_P^{-1}Q\gamma_P$ and
$g_P=\gamma_P^{-1}g'\gamma_P$. The finite-cover equidistribution clause in
D, applied to the common finite-index lattice, makes $g_P$ totally
equidistributed in $Q_P/(Q_P\cap\Gamma)$ on every retained dense progression.
Rationality of $\gamma_P$ gives rationality of $Q_P$.

Partition into those residue progressions and then into intervals of length
$\delta N$ for a sufficiently small standard $\delta$. Averaging selects a
piece with an average bounded below by the original fixed threshold.
On that piece freeze $\varepsilon(n)\gamma_P$ at one point. The uniform
Lipschitz error is $O(\delta)$, so equidistribution gives
\begin{equation}\label{eq:coset}
 \left|\int_{P/(P\cap\Gamma)}F(\varepsilon_0x)\,d\mu(x)\right|>c>0
\end{equation}
for the connected rational $P=Q_P$ and bounded $\varepsilon_0$.
This also holds for an internal test with a standard Lipschitz bound:
compactness gives a finite standard uniform net in that Lipschitz ball.
Apply total equidistribution to each net entry and then let the fixed net
error tend to zero. Thus no unjustified standardness of the selected $F_h$
or of the frozen bounded translate is being used.

We claim $\pi_i(P\cap G_i)\supseteq\Xi_i^\perp$. Let $\lambda$ be an
integer character annihilating the rational space $\pi_i(P\cap G_i)$.
Then $\lambda(\Tay_i(g_P))=0$. Conjugation changes no horizontal
coefficient since $[G,G_i]\subseteq G_{(i,2)}$. Hence the horizontal
homomorphism identity gives
\[
 \lambda(\Tay_i(g))=
 \lambda(\Tay_i(\varepsilon))+
 \lambda(\Tay_i(\gamma)) \pmod1.
\]
The first term is $O(N^{-i})$. The second is rational of bounded standard
denominator: derivatives of $\gamma$ lie in $\widehat\Gamma\cap G_i$, whose
horizontal image is a finite-index over-lattice of that of $\Gamma\cap G_i$.
Thus for some standard positive $q$, $q\lambda\in\Xi_i$. Every annihilator
of the rational image has this property. Taking real annihilators proves the
claim. In particular no false assertion that $\gamma$ takes values in
$\Gamma$ or has zero horizontal coefficient is needed.

\section{Applying the coset conclusion to the actual quadruple}
After all the preceding dense restrictions, reapply hereditary LP03.
There are $\gg N^3$ biased quadruples. The exceptional sets for frequency
independence have size $o(N^3)$. Permutations preserving the additive
quadruple equation make the omitted linear-petal position any of the four;
the union of the four exceptional sets is still sparse. Fix a biased
quadruple with all four required independence statements.

Use the representations of its four nilcharacters and the lower bias test
to form a product group
\[
 \mathcal G=G^{(1)}\times G^{(2)}\times G^{(3)}\times G^{(4)}\times G^{(0)}.
\]
The first four groups have degree--rank at most $(m,R)$; the last has degree
at most $(m,R-1)$, including all four lower representation models and the
ordinary degree-at-most-$m-1$ bias witness. The stronger strict inequality
printed for this last group is unnecessary and is not assumed.
Translations in $n$ preserve horizontal Taylor coefficients by their
homomorphism/difference identity. Zero extensions in \eqref{eq:quad}
introduce only an interval-intersection mask. Approximate this mask in
normalized $L^1$ to error $\delta$ by a bounded Lipschitz circle test in
$n/(cN)$, with fixed large $c$ and Lipschitz constant $O(1/\delta)$.
Its two transition strips have total length $O(\delta N)$; multiplying by
the uniformly bounded remaining tests controls the average error. Choosing
$\delta$ below the fixed bias threshold preserves positive bias. This circle
has degree one and belongs in $G^{(0)}$ because $m\geq2$.

The resulting product test has a nonzero scalar average and exact vertical law
\[
 F(\phi_1(t_1)x_1,\ldots,\phi_4(t_4)x_4,x_0)
 =\ee(\eta(t_1)+\eta(t_2)-\eta(t_3)-\eta(t_4))F(x_1,\ldots,x_4,x_0).
\]
All factors, including the cutoff, are uniformly Lipschitz on a fixed
standard product geometry. Apply the preceding factorization consequence
to obtain a connected rational $P\subseteq\mathcal G$ and the nonzero
coset integral \eqref{eq:coset}.

\section{Rational duality and the two commutator conclusions}
For each weight $i$, let $V_{123,i}$ be the real space spanned by the
three-coordinate diagonal images of core letters and by each individual
coordinate image of every petal letter in positions $1,2,3$. Define
$V_{124,i}$ likewise. The projection of $\pi_i(P\cap\mathcal G_i)$ to
these three factors contains the indicated $V$.

Here is the exact duality argument. If containment failed, rationality gives
a rational linear functional that vanishes on the projected space but not
on $V$. Extend it by zero in the omitted factors and clear denominators to
make it an integer character $\lambda$. Because the full space contains
$\Xi_i^\perp$, this character is in the rational span of $\Xi_i$. After
multiplication by another standard integer it belongs to $\Xi_i$ itself.
Evaluating it on the Taylor coefficient gives a relation
\[
 \sum_{a\in C}(c_{1a}+c_{2a}+c_{3a})\xi_{C,a}
       +\sum_{k=1}^3\sum_{a\in L\cup B}c_{ka}\xi_{h_k,a}
       =O(N^{-i}).
\]
The selected independence forces every petal coefficient and each displayed
core sum to vanish. That says the character vanishes on $V$, a contradiction.
The $124$ proof is identical. This uses rational real subspaces, not an
unjustified identification of an arbitrary real annihilator with one
integral member of $\Xi_i$.

Similarly the projection onto all four main factors contains the space
$V_{B,i}$ generated by diagonal core images and by the individual-coordinate
images of every letter in $B$. In the analogous integer relation eliminate
the fourth $L$ coefficient using
$\xi_{h_4,a}=\xi_{h_1,a}+\xi_{h_2,a}-\xi_{h_3,a}$ for $a\in L$.
Independence makes the four $B$ coefficients vanish separately and the sum
of core coefficients vanish. This is exactly annihilation of $V_{B,i}$.
No numerical range for $B$ is borrowed from the conflicting printed order
of the $L$ and $B$ blocks.

Let $w$ be any bracketing of $R$ letters of weights summing to $m$.
First suppose it has two petal occurrences. For one occurrence lift its
first-coordinate image, with zero second and third coordinates, from
$V_{123,i}$ to $P\cap\mathcal G_i$. For the second lift its first-coordinate
image with zero second and fourth coordinates from $V_{124,i}$. Lift the
remaining letters so that their first horizontal coordinate is the required
one. Take the same iterated commutator of these lifts.

In the first factor it equals $\phi_1(w)$ exactly. Replacing one lift by its
horizontal error inserts an element of $G_{(i,2)}$ in place of
$G_{(i,1)}$, raising the total rank to at least $R+1$ at weight $m$;
every such term vanishes. Higher BCH terms have larger weight and vanish
as well. In factor two, either selected petal lift has such an error;
factor three has the first, and factor four the second. Therefore those
coordinates are the identity. In factor zero, an $R$-fold commutator of
the prescribed weights lies in its trivial $(m,R)$ group. We obtain
\[
                 (\phi_1(w),1,1,1,1)\in P.                 \tag{*}
\]
This reasoning applies to each position of a repeated letter and to every
bracketing, not only a single left-nested commutator.

If $w$ has a $B$ letter, use $V_{B,i}$ to lift that occurrence with all
other main coordinates zero horizontally. Lift the remaining occurrences
by their first coordinates and repeat the same rank-error calculation.
This proves (*) as well. It works for $R=1$; the other top coordinates then
already vanish, and the lower model has trivial weight-$m$ group.

Finally (*) gives a real-zero conclusion. The element is central in
$\mathcal G$. Since $P$ is a connected subgroup of a simply connected
nilpotent group, its entire one-parameter subgroup belongs to $P$.
Left translation by its $t$th power leaves Haar measure on $P/(P\cap\Gamma)$
invariant and commutes with the bounded coset translate. The vertical law
multiplies the nonzero integral by $\ee(t\eta(w))$. Hence
$\ee(t\eta(w))=1$ for every real $t$, so $\eta(w)=0$ as a real number.
Using only $t=1$ would give an integer condition and would be insufficient.
This proves both annihilations in Theorem~\ref{thm:output}.

\section{The LP04 and rank-family joins}
The output has the exact LP04 hypotheses: the corrected universal weighted
group, an integral top character with all required real-zero annihilations,
fixed lattice-compatible $\phi$, matched horizontal Taylor tori, and a unit
representation with a lower degree--rank factor. Its petal quotient kills
the whole ideal generated by the bad words at every weight; the vanishing
just proved concerns the top restriction needed to descend the character.
These distinct roles are not identified.

The Section 12 join chooses its stronger integral extension to the ordinary
top before invoking LP04 comparison, and gives one global multidegree
$(1,m)$ nilcharacter equal to a fixed coordinate embedding of $\zeta_h$.
LP04 then supplies the whole matrix
$(\iota\chi_h)\otimes\overline{\zeta_h}$ with strictly lower degree--rank.
Projecting the first tensor factor through the fixed coordinate projection
gives $\chi_h\otimes\overline{\zeta_h}$ with the same bound and exact values.
Thus the unit padding used in rational removal does not weaken the actual
matrix input required by rank-iteration001.

All selections, roots, covers, maps and tests were made internally on
representatives; each retained dimension, character and bound is standard.
There are finitely many weights and coordinates at any one rank, so all
density and correlation losses remain positive standard numbers. The next
rank family from rank-iteration001 again has exactly the hypotheses of
Section 1. Repeating the present reconstruction at that new rank uses its
own new frequency data. No independence is transported through a
Stone--Weierstrass approximation or assumed for the newly selected test.
There are $m$ rank steps, so the same finite-loss conclusion holds across
the full iteration. This is the required source-to-consumer compatibility,
conditional on the two separately supplied analytic contracts and the
independent acceptance of the project proof units used here.
Representation equality is required on $[N]$. Whenever a global sequence
is used, take the displayed polynomial model as its extension. Small-frequency
removal need not preserve a previous extension outside $[N]$; the interval
mask in the quadruple argument is precisely what makes that change harmless.

\section{Bounded correction and consistency record}
The source scope is final GTZ Sections 9--11, the immediate pp.~1307--1309
consumer, the D.4--D.6 interface, B.4/B.7, and the selection statements in
Appendix A. The attached record pins exact passages and the current project
consumers. Concrete changes are:
\begin{enumerate}
\item State the fixed rank $R$ explicitly in Definition 9.1's second
truncation relation, where the printed formula writes $r$ without the fixed
rank subscript. The printed weight cutoff $\sum i_j\geq s$ is already
equivalent to $\sum i_j>m$ and is not an error. Rendering the source rather
than relying on text extraction establishes this distinction.
\item Retain the common division from 10.9 in 10.10 and restore Freiman laws
after real coefficient division. Fixed arbitrary-location boxes suffice.
\item Rational removal uses a finite cover and explicit same-character unit
completion with a fixed isometry, rather than unproved unit interpolation.
\item Keep varying uniformly Lipschitz $F_h$; total boundedness does not make
them literally equal as claimed in the immediate p.~1308 narrative.
\item Use named core/linear/independent blocks, because p.~1302 reverses the
order used in the statement of 10.10 and Theorem 11.1.
\item Treat the repeated clauses (iii),(iv) of Theorem 11.1 as duplicates;
there are two distinct annihilation conclusions.
\item Use degree--rank at most $(m,R-1)$ for the lower product model, not an
unneeded strict inequality below that index.
\item Retain the rational Taylor term in D.5 and clear its denominator in
the annihilator argument. The factorization subgroup may equal the ambient
group. The separate analytic supplier addresses the D.2 total-equidistribution
issue and F.1 signed-graph/no-wrap details.
\item The real-zero character argument uses all real powers in the connected
subgroup. One integral phase identity alone is insufficient.
\item Reapply LP03 after each dense restriction; intersect many biased
quadruples with almost-all independence, not two unrelated dense sets.
\end{enumerate}
These are authored repairs and scoped dependencies, not a declaration that
all historical versions or the entire inverse theorem have been checked.
The original published source and all earlier frozen packets are preserved.
\end{document}
